\documentclass[11pt]{article}
\usepackage{amsmath}
\usepackage{listings}
\usepackage{hyperref}
\usepackage{graphicx}
\usepackage{booktabs}
\usepackage{geometry}
\geometry{margin=1in}

% Code listing configuration
\lstset{
    basicstyle=\ttfamily\small,
    breaklines=true,
    frame=single,
    numbers=left,
    numberstyle=\tiny,
    showstringspaces=false,
}

\title{GGL90 Python Port Description:\\
       Implementation of MITgcm pkg/ggl90 in 1D Column Model}
\author{Python Port Development Team}
\date{August 2026}

\begin{document}
\maketitle
\tableofcontents
\newpage

\section{Introduction and Scope}
\label{sec:introduction}

The GGL90 Python port is a faithful 1D ocean column implementation of the MITgcm's \verb|pkg/ggl90| turbulence closure scheme. This document describes how the Python code implements the Fortran algorithms and documents key translation decisions, bug fixes, and validation milestones.

\subsection{Purpose of the Python Port}

The Python implementation serves three primary purposes:

\begin{enumerate}
    \item \textbf{Educational platform}: Provide a readable, well-documented reference implementation for understanding GGL90 physics and numerical methods
    \item \textbf{Validation testbed}: Enable rapid debugging of coordinate conventions, sign conventions, and algorithmic correctness in a controlled 1D setting
    \item \textbf{Research tool}: Support sensitivity studies, parameter tuning experiments, and comparative analysis with alternative mixing schemes (e.g., KPP)
\end{enumerate}

\subsection{Relationship to MITgcm Fortran}

The Python port is designed as a \textit{faithful translation} of MITgcm's GGL90 implementation, preserving:

\begin{itemize}
    \item Core physics equations (TKE budget, mixing length formulation, stability functions)
    \item Numerical methods (implicit TKE solver, tridiagonal matrix solution)
    \item Parameter defaults (matching MITgcm defaults, with ECCOv4 R4 as an alternative configuration)
    \item Coordinate conventions (z-axis positive upward, depth positive downward)
\end{itemize}

\textbf{Key simplifications} for 1D column context:
\begin{itemize}
    \item No horizontal mixing or advection
    \item No IDEMIX internal wave energy model
    \item No Langmuir circulation parameterization
    \item No grid-scale noise smoothing (not needed in 1D)
    \item Simplified equation of state (JMD95 without full compressibility terms)
\end{itemize}

\subsection{Target Use Cases}

\begin{itemize}
    \item \textbf{Parameter sensitivity studies}: Test impact of GGL90 tuning parameters (e.g., $\alpha$, $c_k$, $c_\varepsilon$) on mixing profiles
    \item \textbf{Scheme comparison}: Quantify differences between GGL90 and KPP under identical forcing scenarios
    \item \textbf{Training data generation}: Produce controlled scenario outputs for machine learning applications
    \item \textbf{Debugging}: Isolate and validate individual physics components (e.g., mixing length, TKE production/dissipation terms)
\end{itemize}

\subsection{Directory Structure Overview}

\begin{verbatim}
GGL90_ML/
  GGL90_PY/
    ggl90_core_driver.py          # Main orchestrator (≈ ggl90_calc.F)
    ggl90_scheme_specific.py      # Mixing length, TKE terms
    ggl90_mixing_coefficients.py  # Viscosity/diffusivity assembly
    ggl90_parameters.py           # YAML parameter loading
main/
  mixing_adapter.py               # Unified interface for GGL90/KPP
  unified_driver.py               # Time-stepping orchestrator
  physics_basis.py                # N², S², EOS shared utilities
configuration_yamls/
  ggl90_default.yaml              # MITgcm default parameters
  ggl90_eccov4r4.yaml             # ECCOv4 R4 configuration
  physical_parameters.yaml        # Shared constants (g, ρ₀, etc.)
simulations/scenarios/
  scenario_*_initial_conditions.yaml
  scenario_*_atmospheric_forcing.yaml
  scenario_*_time_integration.yaml
\end{verbatim}

\section{Development History and Bug Fixes}
\label{sec:history}

The GGL90 Python port was developed in parallel with the KPP port during 2024-2026. Initial development focused on algorithm translation; subsequent iterations addressed critical bugs related to coordinate conventions, sign conventions, and numerical stability. This section documents the major bug fixes implemented in July 2026 that brought the implementation to production readiness.

\subsection{Timeline}

\begin{itemize}
    \item \textbf{2024-2025}: Initial translation from MITgcm Fortran, basic validation with idealized scenarios
    \item \textbf{2026-Q1}: Integration with unified driver architecture, YAML configuration system
    \item \textbf{2026-Q2}: Development of realistic forcing scenarios (hurricane, Arctic convection, etc.)
    \item \textbf{2026-07-16}: Major bug fix release (Fixes 1-8 documented below), full scenario validation
    \item \textbf{2026-08}: Documentation consolidation, production-ready release
\end{itemize}

\subsection{Fix 1: N² Sign Convention}
\label{sec:fix-n2}

\textbf{Issue}: Buoyancy frequency squared ($N^2$) computed with incorrect sign, causing inverted TKE dissipation behavior.

\textbf{Formula error}:
\begin{equation}
N^2 = \frac{g}{\rho_0} \frac{d\rho}{dz} \quad \text{(WRONG)}
\end{equation}
should be:
\begin{equation}
N^2 = -\frac{g}{\rho_0} \frac{d\rho}{dz} \quad \text{(CORRECT)}
\end{equation}

\textbf{Physical impact}: In stably stratified layers ($d\rho/dz < 0$), the wrong sign caused $N^2 < 0$, leading to negative TKE dissipation terms and unphysical TKE growth in stable regions.

\textbf{Code comparison}:

\textit{Before fix (main/physics\_basis.py)}:
\begin{lstlisting}[language=Python, caption={Incorrect N² computation}]
def compute_buoyancy_frequency_squared(
    rho, grid, g=9.81, rho0=1027.5
):
    """Compute N^2 = (g/rho0) * drho/dz (WRONG SIGN)"""
    drho_dz = np.gradient(rho, grid.depth)
    N_squared = (g / rho0) * drho_dz  # Missing minus sign
    return N_squared
\end{lstlisting}

\textit{After fix}:
\begin{lstlisting}[language=Python, caption={Correct N² computation}]
def compute_buoyancy_frequency_squared(
    rho, grid, g=9.81, rho0=1027.5
):
    """Compute N^2 = -(g/rho0) * drho/dz (CORRECT)"""
    drho_dz = np.gradient(rho, grid.depth)
    N_squared = -(g / rho0) * drho_dz  # Correct sign
    return N_squared
\end{lstlisting}

\textbf{MITgcm comparison}: MITgcm's \verb|pkg/ggl90/ggl90_calc.F| correctly implements the minus sign in the buoyancy frequency calculation.

\textbf{Validation}: Unit test reference formulas updated; all 8 regression tests passed after fix. Scenario outputs showed physically reasonable stratification-dependent mixing.

\subsection{Fix 2: Z-Coordinate Double-Negation}
\label{sec:fix-z-coordinate}

\textbf{Issue}: \verb|ColumnGrid.z_positive_up| property incorrectly negated depth values that were already stored in the correct (negative-down) convention.

\textbf{Root cause}: Property assumed depth array stored as positive-downward, but the grid constructor already stored depth as negative-downward (z $<$ 0 below surface).

\textbf{Impact}: All depth-dependent calculations in GGL90 (mixing length to surface, mixing length to bottom) used inverted z-values, producing incorrect mixing length limiters.

\textit{Before fix (main/column\_grid.py)}:
\begin{lstlisting}[language=Python, caption={Incorrect z-coordinate property}]
@property
def z_positive_up(self):
    """Return coordinates with positive up."""
    return -self.depth  # WRONG: double-negates already-correct values
\end{lstlisting}

\textit{After fix}:
\begin{lstlisting}[language=Python, caption={Correct z-coordinate property}]
@property
def z_positive_up(self):
    """Return coordinates with positive up (depth already negative-down)."""
    return self.depth  # CORRECT: no negation needed
\end{lstlisting}

\textbf{Cascading fixes}: The same bug was found in \verb|diagnostics.py|, where depth-dependent diagnostic computations also incorrectly negated z-values.

\subsection{Fix 3: Mixing-Length Depth Calculations}
\label{sec:fix-mixing-length}

\textbf{Issue}: Mixing-length limiter calculations produced negative distances to surface and bottom, rendering the limiter inert.

\textbf{Root cause}: Used incorrect z-coordinate sign convention (from Fix 2 above):
\begin{lstlisting}[language=Python, caption={Incorrect depth distances}]
depth_to_surface = z - z[0]       # WRONG: negative for z < z[0]
depth_to_bottom = z[-1] - z       # WRONG: negative for z > z[-1]
\end{lstlisting}

\textbf{Physical impact}: Mixing length clamped to background floor everywhere, making GGL90 unresponsive to wind forcing. Example: \verb|hurricane_wind| scenario showed no SST cooling (20.8°C final vs.\ 28°C initial), indicating complete failure of wind-driven mixing.

\textit{After fix (GGL90\_ML/GGL90\_PY/ggl90\_scheme\_specific.py)}:
\begin{lstlisting}[language=Python, caption={Correct depth distance calculations}]
# z[0] is surface (most negative), z[-1] is bottom (least negative)
depth_to_surface = z[0] - z  # Positive distance to surface
depth_to_bottom = z - z[-1]  # Positive distance to bottom

# Apply von Karman mixing length limiter
l_vonkarman_surface = params.GGL90_VONK * depth_to_surface
l_vonkarman_bottom = params.GGL90_VONK * depth_to_bottom
mixing_length = np.minimum(
    mixing_length,
    np.minimum(l_vonkarman_surface, l_vonkarman_bottom)
)
\end{lstlisting}

\textbf{Validation}: After fix, \verb|hurricane_wind| GGL90 final SST = 13.5°C, showing appropriate wind-driven entrainment and thermocline erosion.

\subsection{Fix 4: EOS Temperature Clamp}
\label{sec:fix-eos-clamp}

\textbf{Issue}: JMD95 equation of state polynomial exhibits non-monotonic behavior below $\approx -13°$C, where density \textit{decreases} with further cooling (violates physics).

\textbf{Physical problem}: Unclamped EOS allowed runaway cooling in Arctic scenarios:
\begin{itemize}
    \item Colder water becomes lighter (non-physical)
    \item Triggers positive feedback loop (cooling $\to$ lighter surface $\to$ reduced mixing $\to$ more cooling)
    \item \verb|arctic_convection| scenario reached $-53°$C SST
\end{itemize}

\textit{Solution (main/eos.py)}:
\begin{lstlisting}[language=Python, caption={JMD95 EOS with temperature clamp}]
EOS_MIN_THETA_C = -2.0  # Freezing point clamp

def jmd95_eos(theta, salt, pressure):
    """
    JMD95 equation of state with temperature clamp.
    Clamps theta to prevent non-monotonic density behavior below -2°C.
    """
    # Clamp temperature to valid polynomial range
    theta = np.maximum(theta, EOS_MIN_THETA_C)

    # Polynomial coefficients (unchanged from MITgcm)
    ...
    # Compute density using JMD95 polynomial
    rho = ...
    return rho
\end{lstlisting}

\textbf{Note}: Clamp only affects EOS computation (used for buoyancy frequency $N^2$), not prognostic temperature state. Prognostic state can still cool below $-2°$C; density computation saturates.

\textbf{Validation}: \verb|arctic_convection| now stabilizes at $-8.5°$C, physically reasonable for Arctic winter conditions with sea ice heat flux.

\subsection{Fix 5: Pressure Scale (10$\times$ Error)}
\label{sec:fix-pressure}

\textbf{Issue}: Pressure used in EOS calculation was 10$\times$ too large due to incorrect dbar-to-meter conversion assumption.

\textit{Before fix (GGL90\_ML/GGL90\_PY/ggl90\_core\_driver.py)}:
\begin{lstlisting}[language=Python, caption={Incorrect pressure calculation}]
# WRONG: Assumes 10 meters per dbar (incorrect)
pressure = -10.0 * grid.depth
\end{lstlisting}

\textbf{Correct conversion}: 1 dbar $\approx$ 1 meter of seawater (not 10 meters).

\textit{After fix}:
\begin{lstlisting}[language=Python, caption={Correct pressure calculation}]
# CORRECT: 1 dbar ≈ 1 m seawater
pressure = -grid.depth  # depth is negative, pressure is positive
\end{lstlisting}

\textbf{Physical impact}: 10$\times$ overpressure caused:
\begin{itemize}
    \item In-situ density 10$\times$ too heavy
    \item Overly large $N^2$ (stratification overestimated)
    \item Weak diffusivity (dissipation term too strong)
    \item Systematic underestimation of deep mixing
\end{itemize}

\textbf{Validation}: Full scenario suite re-run; results now show sensible deep mixing magnitudes consistent with KPP expectations.

\subsection{Fix 6: Convective Adjustment (ivdc\_kappa)}
\label{sec:fix-ivdc}

\textbf{Issue}: MITgcm's convective adjustment feature (\verb|ivdc_kappa|) not implemented in Python port, causing disagreement in unstable-column scenarios.

\textbf{MITgcm approach}: In statically unstable regions ($N^2 < 0$), apply enhanced diffusivity:
\begin{equation}
\kappa_{\text{convective}} = \max(\kappa_{\text{GGL90}}, \text{ivdc\_kappa})
\end{equation}

\textit{Implementation (main/unified\_driver.py)}:
\begin{lstlisting}[language=Python, caption={Convective adjustment implementation}]
def _apply_convective_adjustment(self, kappa_e, t_old):
    """
    Apply convective adjustment (ivdc_kappa) to unstable regions.

    Mimics MITgcm behavior: where N^2 < 0, set diffusivity to
    max(computed_kappa, ivdc_kappa).
    """
    # Compute density from old temperature/salinity state
    rho_old = compute_density(t_old, self.s, grid, params)

    # Identify static instability
    N_squared = compute_buoyancy_frequency_squared(
        rho_old, self.grid, g=params.g, rho0=params.rho0
    )
    instability_mask = N_squared < 0.0

    # Apply enhanced diffusivity where unstable
    kappa_e[instability_mask] = np.maximum(
        kappa_e[instability_mask],
        self.ivdc_kappa
    )
    return kappa_e
\end{lstlisting}

\textbf{Configuration}: Added \verb|ivdc_kappa| parameter to \verb|physical_parameters.yaml| (default 0.0 m²/s, ECCOv4 R4 uses 10.0 m²/s).

\textbf{CLI integration}: Both execution scripts support \verb|--ivdc-kappa <value>| flag for runtime override.

\subsection{Fix 7: run\_experiment\_example.py Default Parameters}

\textbf{Issue}: Script hardcoded \verb|ggl90_realistic.yaml| as default, preventing users from easily testing MITgcm default parameters.

\textit{Before fix}:
\begin{lstlisting}[language=Python, caption={Hardcoded parameter file}]
ggl90_params = GGL90Parameters.from_yaml(
    "../configuration_yamls/ggl90_realistic.yaml"
)
\end{lstlisting}

\textit{After fix}:
\begin{lstlisting}[language=Python, caption={Optional parameter override}]
ggl90_params = GGL90Parameters.from_yaml(ggl90_yaml)
# where ggl90_yaml defaults to None, loading built-in MITgcm defaults
\end{lstlisting}

\textbf{New CLI features}:
\begin{itemize}
    \item \verb|--ggl90-yaml <path>|: Override GGL90 parameters
    \item \verb|--kpp-yaml <path>|: Override KPP parameters
    \item \verb|--ivdc-kappa <value>|: Override convective adjustment parameter
    \item \verb|--no-plots|: Suppress plot generation
\end{itemize}

\subsection{Fix 8: Realistic Initial Conditions}

\textbf{Issue}: Scenarios used uniform temperature profiles (e.g., 22°C throughout column), unrealistic for validating mixing schemes.

\textbf{Improvements}: Updated initial condition YAML files with realistic stratification:

\textit{hurricane\_wind scenario}:
\begin{itemize}
    \item Surface mixed layer: 28°C (warm tropical water)
    \item Sharp thermocline: 40-200 m depth (28°C $\to$ 8°C)
    \item Deep water: 8°C $\to$ 5°C
    \item Salinity: 35.2 psu surface, gradual decrease with depth
\end{itemize}

\textit{combined\_storm scenario}:
\begin{itemize}
    \item North Atlantic profile (40-50°N)
    \item Surface: 15°C
    \item Thermocline: 30-200 m (15°C $\to$ 5°C)
    \item Deep: 5°C $\to$ 4°C
    \item Salinity: 35.4 psu (Atlantic water)
\end{itemize}

\textbf{Validation outcomes}: After realistic ICs, scenario final SSTs physically reasonable (GGL90: 13.5°C, KPP: 9.2°C for \verb|hurricane_wind|).

\subsection{Validation Status}

\textbf{Regression tests}: 8/8 passing (test\_staggering.py)

\textbf{Scenario suite}: All 6 scenarios $\times$ 2 schemes validated, exit code 0

\begin{table}[h]
\centering
\caption{Scenario Validation Results (2026-07-16)}
\begin{tabular}{lcc}
\toprule
Scenario & KPP SST (°C) & GGL90 SST (°C) \\
\midrule
arctic\_convection & $-5.99$ & $-8.49$ \\
calm\_baseline & 21.89 & 21.91 \\
combined\_storm & 7.21 & 8.27 \\
heavy\_rain\_freshening & 21.97 & 21.96 \\
hurricane\_wind & 9.23 & 13.49 \\
tropical\_heating\_diurnal & 26.36 & 26.39 \\
\bottomrule
\end{tabular}
\end{table}

\section{File Organization and Module Structure}
\label{sec:file-organization}

The GGL90 Python implementation follows a modular architecture separating scheme-specific physics (GGL90 core), shared utilities (EOS, grid, solvers), and unified driver infrastructure.

\subsection{Core GGL90 Modules}

\textbf{ggl90\_core\_driver.py} (main orchestrator, $\approx$ 450 lines)

Python equivalent of MITgcm's \verb|ggl90_calc.F|. Orchestrates the complete GGL90 calculation sequence:
\begin{enumerate}
    \item Stratification and shear diagnosis ($N^2$, $S^2$) via \verb|physics_basis.py|
    \item Mixing length computation via \verb|ggl90_scheme_specific.py|
    \item TKE production/buoyancy/dissipation terms
    \item Implicit TKE solver (tridiagonal matrix)
    \item Viscosity/diffusivity assembly via \verb|ggl90_mixing_coefficients.py|
\end{enumerate}

Key class: \verb|GGL90Driver|, with method \verb|compute_mixing()| as main entry point.

\textbf{ggl90\_scheme\_specific.py} (GGL90-specific physics, $\approx$ 300 lines)

Implements GGL90-specific formulations:
\begin{itemize}
    \item \verb|GGL90MixingLength|: Mixing length limiter (von Kármán, min/max bounds)
    \item \verb|compute_tke_production()|: Shear production term $P = \kappa_m S^2$
    \item \verb|compute_tke_buoyancy()|: Buoyancy term $B = -\kappa_h N^2$
    \item \verb|compute_tke_dissipation()|: Dissipation term $\varepsilon = c_\varepsilon k^{3/2} / \ell$
\end{itemize}

Corresponds to physics computations in \verb|ggl90_calc.F| lines 200-400 (MITgcm).

\textbf{ggl90\_mixing\_coefficients.py} (diffusivity assembly, $\approx$ 150 lines)

Computes viscosity and diffusivity from TKE and mixing length using stability functions:
\begin{align}
\kappa_m &= c_k \sqrt{k} \ell S_m \\
\kappa_h &= c_k \sqrt{k} \ell S_h
\end{align}
where $S_m$, $S_h$ are stability functions depending on local Richardson number.

Corresponds to \verb|ggl90_calc_visc.F| and \verb|ggl90_calc_diff.F| in MITgcm.

\textbf{ggl90\_parameters.py} (configuration, $\approx$ 200 lines)

Parameter dataclass with YAML loader:
\begin{lstlisting}[language=Python, caption={GGL90Parameters class}]
@dataclass
class GGL90Parameters:
    """GGL90 parameter configuration (matches MITgcm defaults)."""
    GGL90ck: float = 0.1
    GGL90ceps: float = 0.7
    GGL90alpha: float = 1.0  # ECCOv4 R4: 30.0
    GGL90m2: float = 3.75
    GGL90TKEmin: float = 1e-11
    GGL90TKEsurfMin: float = 1e-4
    GGL90TKEbottom: float = 1e-11
    GGL90viscMax: float = 100.0
    GGL90diffMax: float = 100.0
    mxlMaxFlag: int = 0
    mxlSurfFlag: bool = False

    @classmethod
    def from_yaml(cls, yaml_path=None):
        """Load from YAML file, or use defaults if None."""
        ...
\end{lstlisting}

\subsection{Shared Utilities (main/ directory)}

\textbf{physics\_basis.py} (stratification/shear, $\approx$ 200 lines)

Shared physics computations used by both GGL90 and KPP:
\begin{itemize}
    \item \verb|compute_buoyancy_frequency_squared()|: $N^2 = -(g/\rho_0) d\rho/dz$
    \item \verb|compute_vertical_shear_squared()|: $S^2 = (du/dz)^2 + (dv/dz)^2$
\end{itemize}

\textbf{eos.py} (equation of state, $\approx$ 150 lines)

JMD95 equation of state implementation with temperature clamp (Fix 4).

\textbf{shared\_column\_solver.py} (tridiagonal solver, $\approx$ 100 lines)

Thomas algorithm for tridiagonal systems arising in implicit diffusion.

\textbf{column\_grid.py} (vertical grid, $\approx$ 150 lines)

Vertical grid specification with staggering conventions:
\begin{lstlisting}[language=Python, caption={ColumnGrid class}]
class ColumnGrid:
    """
    1D vertical grid with MITgcm-compatible staggering.

    - depth: cell centers (negative-downward, z < 0)
    - z_positive_up: cell centers (positive-upward, z > 0)
    - depth_f: cell interfaces
    """
    def __init__(self, depth):
        self.depth = depth
        self.nz = len(depth)
        ...

    @property
    def z_positive_up(self):
        return self.depth  # No negation (Fix 2)
\end{lstlisting}

\subsection{Unified Driver Infrastructure}

\textbf{mixing\_adapter.py} (scheme abstraction, $\approx$ 250 lines)

Provides unified interface for GGL90 and KPP, abstracting differences:
\begin{lstlisting}[language=Python, caption={MixingAdapter interface}]
class MixingAdapter:
    """Unified interface for mixing schemes."""

    def __init__(self, scheme: str, params):
        if scheme == "ggl90":
            self.driver = GGL90Driver(params)
        elif scheme == "kpp":
            self.driver = KPPDriver(params)

    def compute_mixing(self, state, grid, forcing):
        """
        Compute mixing coefficients.
        Returns: kappa_m, kappa_h, diagnostics
        """
        ...
\end{lstlisting}

\textbf{unified\_driver.py} (time-stepping, $\approx$ 600 lines)

Time-stepping orchestrator integrating mixing computation with advection-diffusion solver:
\begin{enumerate}
    \item Call mixing scheme: \verb|compute_mixing()|
    \item Apply convective adjustment (Fix 6)
    \item Implicit diffusion solve (temperature, salinity)
    \item Apply atmospheric forcing (wind stress, heat flux)
    \item Save diagnostics
\end{enumerate}

\subsection{Configuration Files}

\textbf{configuration\_yamls/ggl90\_default.yaml}

MITgcm default parameters (see \verb|ggl90_parameters.py| for values).

\textbf{configuration\_yamls/ggl90\_eccov4r4.yaml}

ECCOv4 Release 4 tuned parameters:
\begin{verbatim}
GGL90alpha: 30.0        # (default: 1.0)
GGL90TKEmin: 1.0e-7     # (default: 1.0e-11)
GGL90TKEbottom: 1.0e-6  # (default: 1.0e-11)
mxlMaxFlag: 2           # (default: 0)
mxlSurfFlag: true       # (default: false)
\end{verbatim}

\textbf{configuration\_yamls/physical\_parameters.yaml}

Shared physical constants (single source of truth):
\begin{verbatim}
g: 9.81
rho0: 1027.5
cp: 3994.0
ivdc_kappa: 0.0  # (ECCOv4 R4: 10.0)
\end{verbatim}

\subsection{Mapping to MITgcm Files}

\begin{table}[h]
\centering
\caption{Python $\leftrightarrow$ MITgcm File Correspondence}
\begin{tabular}{ll}
\toprule
\textbf{Python File} & \textbf{MITgcm Fortran File} \\
\midrule
ggl90\_core\_driver.py & ggl90\_calc.F \\
ggl90\_scheme\_specific.py & ggl90\_calc.F (physics sections) \\
ggl90\_mixing\_coefficients.py & ggl90\_calc\_visc.F, ggl90\_calc\_diff.F \\
ggl90\_parameters.py & GGL90\_PARMS.h, ggl90\_readparms.F \\
physics\_basis.py & calc\_grad\_phi\_hyd.F, calc\_viscosity.F \\
shared\_column\_solver.py & solve\_tridiagonal.F \\
\bottomrule
\end{tabular}
\end{table}

\section{Coordinate System and Sign Conventions}
\label{sec:coordinates}

Correct coordinate conventions are critical for vertical ocean physics. This section documents the Python port's coordinate system and compares to MITgcm conventions.

\subsection{Vertical Coordinate System}

\textbf{Depth convention}: Depth is stored as negative-downward:
\begin{itemize}
    \item Surface: $z = 0$ (shallowest, most negative in array: \verb|depth[0]|)
    \item Bottom: $z = -H$ (deepest, least negative in array: \verb|depth[-1]|)
\end{itemize}

Example 100-meter column:
\begin{verbatim}
depth = [0, -10, -20, ..., -90, -100]  (negative-downward)
\end{verbatim}

\textbf{Z-positive-up convention} (used in physics calculations):
\begin{itemize}
    \item Surface: $z = 0$
    \item Below surface: $z < 0$
    \item After Fix 2, \verb|z_positive_up| property correctly returns \verb|depth| without negation
\end{itemize}

\subsection{Sign Conventions for Gradients}

\textbf{Vertical derivative sign}: Using $z$ positive upward, derivatives follow standard calculus:
\begin{align}
\frac{d\rho}{dz} &< 0 \quad \text{(stable: density increases downward)} \\
N^2 &= -\frac{g}{\rho_0} \frac{d\rho}{dz} > 0 \quad \text{(stable stratification)}
\end{align}

\textbf{Critical minus sign}: The minus sign in $N^2$ definition is \textit{not optional} (Fix 1). Without it, stable stratification produces negative $N^2$, causing sign errors in TKE dissipation.

\subsection{Flux and Diffusivity Sign Conventions}

\textbf{Diffusive flux convention}:
\begin{equation}
F = -\kappa \frac{d\theta}{dz}
\end{equation}

Upward flux ($F > 0$, warming the layer above) requires $d\theta/dz < 0$ (cold above, warm below) with positive diffusivity $\kappa > 0$.

\textbf{Staggering}: Diffusivity $\kappa$ defined at cell interfaces; tracer values at cell centers. Flux at interface $k+1/2$:
\begin{equation}
F_{k+1/2} = -\kappa_{k+1/2} \frac{\theta_{k+1} - \theta_k}{\Delta z}
\end{equation}

\subsection{Comparison to MITgcm Conventions}

MITgcm uses \textit{r-coordinate} for vertical (positive upward in default configuration):
\begin{itemize}
    \item \verb|rC|: cell centers (positive upward)
    \item \verb|rF|: cell interfaces (positive upward)
    \item Diffusivity arrays staggered at interfaces
\end{itemize}

Python port uses \textit{depth} (negative-downward for storage, positive-upward for physics). Translation between conventions handled by \verb|z_positive_up| property.

\textbf{Key equivalence}:
\begin{verbatim}
MITgcm rC[k]  <->  Python depth[k] (after Fix 2)
MITgcm drF[k] <->  Python grid.dz[k]
\end{verbatim}

\subsection{Distance Calculations (Fix 3)}

Positive distances to boundaries (used in mixing length limiters):
\begin{align}
d_{\text{surface}} &= z_{\text{surface}} - z = 0 - z = -z \quad (z < 0) \\
d_{\text{bottom}} &= z - z_{\text{bottom}} = z - (-H) = z + H
\end{align}

\textit{Implementation}:
\begin{lstlisting}[language=Python, caption={Positive distance computation}]
# z[0] = surface (most negative), z[-1] = bottom (least negative)
depth_to_surface = z[0] - z  # e.g., 0 - (-50) = 50 m
depth_to_bottom = z - z[-1]  # e.g., -50 - (-100) = 50 m
\end{lstlisting}

Before Fix 3, these were computed as \verb|z - z[0]| and \verb|z[-1] - z|, producing negative distances and disabling the mixing length limiter.

\section{Core Algorithm Implementation}
\label{sec:core-algorithm}

This section describes how the Python port implements the GGL90 turbulence closure equations from Gaspar et al. (1990), with comparison to MITgcm's \verb|ggl90_calc.F|.

\subsection{GGL90 Physics Overview}

GGL90 is a prognostic turbulence scheme evolving turbulent kinetic energy (TKE) as a state variable. The TKE budget equation:
\begin{equation}
\frac{\partial k}{\partial t} = P + B - \varepsilon + \frac{\partial}{\partial z}\left(K_k \frac{\partial k}{\partial z}\right)
\end{equation}
where:
\begin{align}
P &= \kappa_m S^2 \quad \text{(shear production)} \\
B &= -\kappa_h N^2 \quad \text{(buoyancy term, positive for convection)} \\
\varepsilon &= c_\varepsilon \frac{k^{3/2}}{\ell} \quad \text{(dissipation)}
\end{align}

Mixing coefficients computed from TKE and mixing length:
\begin{align}
\kappa_m &= c_k \sqrt{k} \ell \, S_m \quad \text{(eddy viscosity)} \\
\kappa_h &= c_k \sqrt{k} \ell \, S_h \quad \text{(eddy diffusivity)}
\end{align}
where $S_m$, $S_h$ are stability functions depending on local Richardson number $Ri = N^2/S^2$.

\subsection{Main Driver: compute\_mixing()}

Entry point for GGL90 calculation (GGL90Driver.compute\_mixing()):

\begin{lstlisting}[language=Python, caption={GGL90 main driver structure}]
def compute_mixing(self, state, grid, forcing, dt):
    """
    Compute GGL90 mixing coefficients for one timestep.

    1. Compute stratification (N²) and shear (S²)
    2. Compute mixing length
    3. Compute TKE source terms (P, B, ε)
    4. Solve implicit TKE equation
    5. Compute viscosity/diffusivity from new TKE
    """
    # Step 1: Stratification and shear
    N_squared = compute_buoyancy_frequency_squared(
        rho, grid, g=params.g, rho0=params.rho0
    )
    S_squared = compute_vertical_shear_squared(
        u, v, grid
    )

    # Step 2: Mixing length
    mixing_length = GGL90MixingLength.compute(
        tke_old, N_squared, grid, params
    )

    # Step 3: TKE source terms
    production = compute_tke_production(
        tke_old, S_squared, mixing_length, params
    )
    buoyancy = compute_tke_buoyancy(
        tke_old, N_squared, mixing_length, params
    )
    dissipation = compute_tke_dissipation(
        tke_old, mixing_length, params
    )

    # Step 4: Implicit TKE solve
    tke_new = self._solve_tke_implicit(
        tke_old, production, buoyancy, dissipation, dt, grid
    )

    # Step 5: Mixing coefficients
    kappa_m, kappa_h = compute_viscosity_diffusivity(
        tke_new, mixing_length, N_squared, S_squared, params
    )

    return GGL90Output(
        tke_new=tke_new, kappa_m=kappa_m, kappa_h=kappa_h,
        mixing_length=mixing_length, n_square=N_squared,
        shear_square=S_squared, ...
    )
\end{lstlisting}

\subsection{Mixing Length Computation}

Mixing length computed as geometric mean of vertical and horizontal scales, with von Kármán limiters near boundaries:

\begin{lstlisting}[language=Python, caption={Mixing length computation}]
@staticmethod
def compute(tke, N_squared, grid, params):
    """
    Compute GGL90 mixing length.

    Following Gaspar et al. (1990) and MITgcm ggl90_calc.F.
    """
    z = grid.z_positive_up
    dz = grid.dz

    # Base mixing length from TKE and stratification
    # ell ~ sqrt(2*k) / N (stratified limit)
    N_safe = np.maximum(N_squared, 1e-14)**0.5
    ell_base = params.GGL90alpha * np.sqrt(2.0 * tke) / N_safe

    # Apply von Karman limiters near boundaries (Fix 3)
    depth_to_surface = z[0] - z  # Positive distance to surface
    depth_to_bottom = z - z[-1]  # Positive distance to bottom

    ell_surface = params.GGL90_VONK * depth_to_surface
    ell_bottom = params.GGL90_VONK * depth_to_bottom

    mixing_length = np.minimum(ell_base, ell_surface)
    mixing_length = np.minimum(mixing_length, ell_bottom)

    # Apply min/max bounds
    mixing_length = np.maximum(mixing_length, params.GGL90mixingLengthMin)
    if params.mxlMaxFlag > 0:
        ell_max = ...  # Scheme-specific maximum
        mixing_length = np.minimum(mixing_length, ell_max)

    return mixing_length
\end{lstlisting}

\textbf{MITgcm comparison}: Corresponds to \verb|ggl90_calc.F| lines 250-320, where mixing length computed with identical von Kármán limiter logic.

\subsection{TKE Source Terms}

\textbf{Shear production}:
\begin{lstlisting}[language=Python, caption={TKE shear production}]
def compute_tke_production(tke, S_squared, mixing_length, params):
    """
    Shear production: P = kappa_m * S²
    where kappa_m = c_k * sqrt(k) * ell * S_m
    """
    sqrt_tke = np.sqrt(np.maximum(tke, params.GGL90TKEmin))

    # Stability function S_m (simplified: constant)
    S_m = params.GGL90m2  # Default: 3.75

    kappa_m = params.GGL90ck * sqrt_tke * mixing_length * S_m

    production = kappa_m * S_squared
    return production
\end{lstlisting}

\textbf{Buoyancy term}:
\begin{lstlisting}[language=Python, caption={TKE buoyancy term}]
def compute_tke_buoyancy(tke, N_squared, mixing_length, params):
    """
    Buoyancy term: B = -kappa_h * N²
    (Positive for convection, negative for stable stratification)
    """
    sqrt_tke = np.sqrt(np.maximum(tke, params.GGL90TKEmin))

    # Stability function S_h (simplified: S_h ~ S_m / Prandtl)
    S_h = params.GGL90m2 / params.GGL90_PRANDTL

    kappa_h = params.GGL90ck * sqrt_tke * mixing_length * S_h

    buoyancy = -kappa_h * N_squared  # Note: N² sign fixed (Fix 1)
    return buoyancy
\end{lstlisting}

\textbf{Dissipation}:
\begin{lstlisting}[language=Python, caption={TKE dissipation}]
def compute_tke_dissipation(tke, mixing_length, params):
    """
    Dissipation: ε = c_ε * k^(3/2) / ell
    """
    sqrt_tke = np.sqrt(np.maximum(tke, params.GGL90TKEmin))
    k_32 = tke * sqrt_tke

    # Mixing length with floor to prevent division by zero
    ell_safe = np.maximum(mixing_length, params.GGL90mixingLengthMin)

    dissipation = params.GGL90ceps * k_32 / ell_safe
    return dissipation
\end{lstlisting}

\subsection{Implicit TKE Solver}

TKE evolved implicitly to maintain numerical stability:
\begin{equation}
\frac{k^{n+1} - k^n}{\Delta t} = P + B - \varepsilon + \frac{\partial}{\partial z}\left(K_k \frac{\partial k^{n+1}}{\partial z}\right)
\end{equation}

Discretized as tridiagonal system:
\begin{equation}
a_i k^{n+1}_{i-1} + b_i k^{n+1}_i + c_i k^{n+1}_{i+1} = d_i
\end{equation}

\begin{lstlisting}[language=Python, caption={Implicit TKE solver}]
def _solve_tke_implicit(self, tke_old, production, buoyancy,
                        dissipation, dt, grid):
    """
    Solve implicit TKE equation using tridiagonal solver.
    """
    nz = grid.nz
    dz = grid.dz

    # Vertical diffusivity for TKE (K_k ~ kappa_m)
    K_k = self._compute_tke_diffusivity(tke_old, mixing_length)

    # Build tridiagonal system: [I - dt*A] k^{n+1} = k^n + dt*S
    # where A is diffusion operator, S is source term
    a = np.zeros(nz)  # Lower diagonal
    b = np.zeros(nz)  # Main diagonal
    c = np.zeros(nz)  # Upper diagonal
    d = np.zeros(nz)  # Right-hand side

    for i in range(1, nz-1):
        # Diffusion coefficients at interfaces
        K_upper = 0.5 * (K_k[i-1] + K_k[i])
        K_lower = 0.5 * (K_k[i] + K_k[i+1])

        # Tridiagonal coefficients
        a[i] = -dt * K_upper / (dz[i] * dz[i-1])
        c[i] = -dt * K_lower / (dz[i] * dz[i+1])
        b[i] = 1.0 - a[i] - c[i]

        # Right-hand side: k^n + dt*(P + B - ε)
        source = production[i] + buoyancy[i] - dissipation[i]
        d[i] = tke_old[i] + dt * source

    # Boundary conditions
    b[0] = 1.0
    d[0] = params.GGL90TKEsurfMin  # Surface TKE minimum
    b[-1] = 1.0
    d[-1] = params.GGL90TKEbottom  # Bottom TKE

    # Solve tridiagonal system
    tke_new = solve_tridiagonal(a, b, c, d)

    # Enforce minimum TKE
    tke_new = np.maximum(tke_new, params.GGL90TKEmin)

    return tke_new
\end{lstlisting}

\textbf{MITgcm comparison}: Identical structure to \verb|ggl90_calc.F| implicit solver (lines 400-500). Python uses Thomas algorithm (\verb|solve_tridiagonal|), MITgcm uses similar direct solver.

\subsection{Viscosity and Diffusivity Assembly}

Final step: compute mixing coefficients from updated TKE:

\begin{lstlisting}[language=Python, caption={Mixing coefficient assembly}]
def compute_viscosity_diffusivity(tke, mixing_length, N_squared,
                                   S_squared, params):
    """
    Compute eddy viscosity and diffusivity from TKE and mixing length.
    """
    sqrt_tke = np.sqrt(np.maximum(tke, params.GGL90TKEmin))

    # Stability functions (simplified: constant)
    S_m = params.GGL90m2
    S_h = S_m / params.GGL90_PRANDTL

    # Mixing coefficients
    kappa_m = params.GGL90ck * sqrt_tke * mixing_length * S_m
    kappa_h = params.GGL90ck * sqrt_tke * mixing_length * S_h

    # Apply maximum bounds
    kappa_m = np.minimum(kappa_m, params.GGL90viscMax)
    kappa_h = np.minimum(kappa_h, params.GGL90diffMax)

    return kappa_m, kappa_h
\end{lstlisting}

\section{Numerical Methods and Time-Stepping}
\label{sec:numerical-methods}

\subsection{Implicit TKE Diffusion}

The vertical TKE diffusion term is discretized implicitly to ensure numerical stability for large diffusivities and coarse vertical grids.

\textbf{Time discretization}:
\begin{equation}
\frac{k^{n+1} - k^n}{\Delta t} = \text{sources} + \frac{\partial}{\partial z}\left(K_k \frac{\partial k^{n+1}}{\partial z}\right)
\end{equation}

The implicit treatment of the diffusion term (right side evaluated at time level $n+1$) removes timestep restrictions from vertical diffusion CFL condition.

\textbf{Spatial discretization}:
\begin{equation}
\left(\frac{\partial}{\partial z} K_k \frac{\partial k}{\partial z}\right)_i \approx \frac{1}{\Delta z_i} \left[ K_{k,i+1/2} \frac{k_{i+1} - k_i}{\Delta z_{i+1/2}} - K_{k,i-1/2} \frac{k_i - k_{i-1}}{\Delta z_{i-1/2}} \right]
\end{equation}

This leads to tridiagonal system solved by Thomas algorithm (\verb|solve_tridiagonal|).

\subsection{Comparison to MITgcm Approach}

MITgcm uses identical implicit approach for TKE equation. Key correspondences:

\begin{table}[h]
\centering
\caption{Numerical Method Comparison}
\begin{tabular}{lll}
\toprule
\textbf{Aspect} & \textbf{MITgcm} & \textbf{Python Port} \\
\midrule
TKE time-stepping & Implicit & Implicit \\
Diffusion discretization & Centered, 2nd-order & Centered, 2nd-order \\
Tridiagonal solver & Direct (custom) & Thomas algorithm \\
Source term treatment & Explicit & Explicit \\
Timestep restriction & CFL from advection & CFL from advection \\
\bottomrule
\end{tabular}
\end{table}

\textbf{Stability}: Implicit diffusion ensures unconditional stability for vertical mixing. Timestep limited only by explicit advection terms (not present in 1D column model) and source term stability.

\subsection{Tracer Diffusion Time-Stepping}

Temperature and salinity evolved using mixing coefficients from GGL90:
\begin{equation}
\frac{\partial \theta}{\partial t} = \frac{\partial}{\partial z}\left(\kappa_h \frac{\partial \theta}{\partial z}\right) + \text{forcing}
\end{equation}

Also treated implicitly (tridiagonal solve in \verb|unified_driver.py|).

\subsection{Timestep Selection}

Typical timestep range: 300-3600 seconds (5 minutes to 1 hour), depending on forcing strength.

\textbf{Stability considerations}:
\begin{itemize}
    \item Implicit diffusion: unconditionally stable
    \item Surface forcing: may require shorter timesteps for rapid heat flux changes
    \item TKE source terms: explicit treatment may require moderate timesteps ($\Delta t \lesssim 3600$ s)
\end{itemize}

Example from \verb|hurricane_wind| scenario: $\Delta t = 1800$ s (30 minutes) maintained stability throughout 5-day simulation.

\section{Parameter Handling and Configuration}
\label{sec:parameters}

\subsection{YAML Configuration System}

The Python port uses YAML files for parameter configuration, replacing MITgcm's namelist (\verb|data.ggl90|) approach.

\textbf{Advantages}:
\begin{itemize}
    \item Human-readable, structured format
    \item Easy to version control and diff
    \item Supports hierarchical parameter organization
    \item Inline comments for documentation
\end{itemize}

\subsection{GGL90 Parameter Sets}

\textbf{Default parameters} (\verb|ggl90_default.yaml|):

Matches MITgcm defaults from \verb|GGL90_PARMS.h|:
\begin{verbatim}
GGL90ck: 0.1           # TKE -> diffusivity coefficient
GGL90ceps: 0.7         # Dissipation coefficient
GGL90alpha: 1.0        # Mixing length coefficient
GGL90m2: 3.75          # Stability function parameter
GGL90TKEmin: 1.0e-11   # Minimum TKE [m²/s²]
GGL90TKEsurfMin: 1.0e-4  # Surface TKE minimum
GGL90TKEbottom: 1.0e-11  # Bottom boundary TKE
GGL90viscMax: 100.0    # Maximum viscosity [m²/s]
GGL90diffMax: 100.0    # Maximum diffusivity [m²/s]
mxlMaxFlag: 0          # Mixing length limiter flag
mxlSurfFlag: false     # Surface mixing length flag
\end{verbatim}

\textbf{ECCOv4 R4 parameters} (\verb|ggl90_eccov4r4.yaml|):

Tuned for global ocean state estimation:
\begin{verbatim}
GGL90alpha: 30.0       # Increased from 1.0
GGL90TKEmin: 1.0e-7    # Increased from 1.0e-11
GGL90TKEbottom: 1.0e-6 # Increased from 1.0e-11
mxlMaxFlag: 2          # Enable maximum mixing length limiter
mxlSurfFlag: true      # Enable surface mixing length limiter
\end{verbatim}

Key differences: ECCOv4 R4 uses larger $\alpha$ (longer mixing lengths), higher background TKE floors, and stricter mixing length limiters.

\subsection{Physical Parameters (Shared)}

Constants shared across schemes in \verb|physical_parameters.yaml|:
\begin{verbatim}
# Fundamental constants
g: 9.81              # Gravitational acceleration [m/s²]
rho0: 1027.5         # Reference density [kg/m³]
cp: 3994.0           # Specific heat capacity [J/(kg·K)]

# Convective adjustment
ivdc_kappa: 0.0      # Convective diffusivity [m²/s]
                      # ECCOv4 R4: 10.0

# EOS parameters
EOS_MIN_THETA_C: -2.0  # Temperature clamp [°C] (Fix 4)
\end{verbatim}

\subsection{Parameter Loading}

\begin{lstlisting}[language=Python, caption={YAML parameter loading}]
@dataclass
class GGL90Parameters:
    """GGL90 parameter dataclass."""
    GGL90ck: float = 0.1
    GGL90ceps: float = 0.7
    # ... (all parameters with defaults)

    @classmethod
    def from_yaml(cls, yaml_path=None):
        """
        Load parameters from YAML file.

        If yaml_path is None, use built-in defaults.
        """
        if yaml_path is None:
            return cls()  # Use defaults

        with open(yaml_path, 'r') as f:
            config = yaml.safe_load(f)

        # Create instance, overriding defaults with YAML values
        return cls(**{k: v for k, v in config.items()
                     if hasattr(cls, k)})
\end{lstlisting}

\subsection{Command-Line Override}

Both execution scripts support parameter override:

\begin{verbatim}
# Use defaults
python run_scenarios.py

# Use ECCOv4 R4 parameters
python run_scenarios.py --ggl90-yaml ../configuration_yamls/ggl90_eccov4r4.yaml

# Override convective adjustment
python run_experiment_example.py --ivdc-kappa 10.0 --ggl90-yaml custom.yaml
\end{verbatim}

\subsection{Parameter Validation}

Basic validation checks applied on parameter load:
\begin{itemize}
    \item Positive definiteness: $c_k, c_\varepsilon, \alpha > 0$
    \item TKE minimums: $\text{TKEmin} > 0$
    \item Diffusivity bounds: $\kappa_{\max} > 0$
\end{itemize}

Invalid parameters raise \verb|ValueError| with descriptive message.

\section{Integration with Model Driver}
\label{sec:integration}

\subsection{Unified Driver Architecture}

The Python implementation uses a unified driver architecture supporting multiple mixing schemes (GGL90, KPP) through a common interface.

\textbf{Architecture components}:
\begin{enumerate}
    \item \textbf{Mixing adapter} (\verb|mixing_adapter.py|): Abstracts scheme-specific interfaces
    \item \textbf{Unified driver} (\verb|unified_driver.py|): Orchestrates time-stepping, diagnostics, I/O
    \item \textbf{Scheme drivers} (\verb|ggl90_core_driver.py|, \verb|kpp_core_driver.py|): Implement scheme-specific physics
\end{enumerate}

\subsection{Mixing Adapter Interface}

\begin{lstlisting}[language=Python, caption={MixingAdapter class}]
class MixingAdapter:
    """
    Unified interface for mixing schemes.

    Provides scheme-agnostic access to mixing computations,
    handling differences in inputs, outputs, and state variables.
    """

    def __init__(self, scheme: str, params):
        """
        Initialize mixing adapter.

        Parameters
        ----------
        scheme : str
            Mixing scheme name ('ggl90' or 'kpp')
        params : GGL90Parameters | KPPParameters
            Scheme-specific parameters
        """
        if scheme == "ggl90":
            self.driver = GGL90Driver(params)
            self.is_prognostic = True  # GGL90 has state (TKE)
        elif scheme == "kpp":
            self.driver = KPPDriver(params)
            self.is_prognostic = False  # KPP is diagnostic
        else:
            raise ValueError(f"Unknown scheme: {scheme}")

    def compute_mixing(self, state, grid, forcing, dt):
        """
        Compute mixing coefficients.

        Returns
        -------
        kappa_m : np.ndarray
            Eddy viscosity [m²/s]
        kappa_h : np.ndarray
            Eddy diffusivity [m²/s]
        diagnostics : dict
            Scheme-specific diagnostics
        """
        output = self.driver.compute_mixing(state, grid, forcing, dt)
        return output.kappa_m, output.kappa_h, output.to_dict()

    def get_state(self):
        """Get prognostic state (TKE for GGL90, None for KPP)."""
        if self.is_prognostic:
            return self.driver.tke
        return None

    def set_state(self, state):
        """Set prognostic state (GGL90 only)."""
        if self.is_prognostic:
            self.driver.tke = state
\end{lstlisting}

\subsection{Time-Stepping Integration}

The unified driver (\verb|UnifiedDriver| class) integrates mixing scheme into time-stepping loop:

\begin{lstlisting}[language=Python, caption={Time-stepping with GGL90 mixing}]
def timestep(self, dt):
    """
    Advance model one timestep.

    1. Compute mixing coefficients (call scheme)
    2. Apply convective adjustment
    3. Diffuse temperature/salinity implicitly
    4. Apply surface forcing
    5. Update diagnostics
    """
    # Step 1: Mixing coefficients from scheme
    kappa_m, kappa_h, diag = self.mixing_adapter.compute_mixing(
        state=self.state,
        grid=self.grid,
        forcing=self.forcing,
        dt=dt
    )

    # Step 2: Convective adjustment (Fix 6)
    if self.ivdc_kappa > 0:
        kappa_h = self._apply_convective_adjustment(
            kappa_h, self.state.temperature
        )

    # Step 3: Implicit diffusion
    self.state.temperature = self._diffuse_tracer(
        self.state.temperature, kappa_h, dt
    )
    self.state.salinity = self._diffuse_tracer(
        self.state.salinity, kappa_h, dt
    )

    # Step 4: Surface forcing
    self._apply_surface_forcing(dt)

    # Step 5: Diagnostics
    self._save_diagnostics(diag)
\end{lstlisting}

\subsection{Scenario Runner}

\verb|run_scenarios.py| provides batch execution capability:

\begin{lstlisting}[language=Python, caption={Scenario runner structure}]
def run_scenarios(scheme, ggl90_yaml=None, kpp_yaml=None,
                  ivdc_kappa=None):
    """
    Run all scenarios for specified scheme(s).

    Parameters
    ----------
    scheme : str
        'ggl90', 'kpp', or 'both'
    ggl90_yaml : Path, optional
        GGL90 parameter override
    kpp_yaml : Path, optional
        KPP parameter override
    ivdc_kappa : float, optional
        Convective adjustment parameter
    """
    scenarios = [
        'arctic_convection',
        'calm_baseline',
        'combined_storm',
        'heavy_rain_freshening',
        'hurricane_wind',
        'tropical_heating_diurnal',
    ]

    results = {}
    for scenario_name in scenarios:
        # Load scenario configuration (IC, forcing, time)
        config = load_scenario_config(scenario_name)

        # Initialize driver
        driver = UnifiedDriver(
            scheme=scheme,
            config=config,
            ggl90_yaml=ggl90_yaml,
            kpp_yaml=kpp_yaml,
            ivdc_kappa=ivdc_kappa,
        )

        # Run simulation
        driver.run()

        # Collect results
        results[scenario_name] = driver.get_summary()

    # Print summary table
    print_summary_table(results)
\end{lstlisting}

\subsection{Diagnostics Integration}

Diagnostics saved at each timestep include:
\begin{itemize}
    \item Mixing coefficients: $\kappa_m$, $\kappa_h$
    \item GGL90-specific: TKE, mixing length, $N^2$, $S^2$, $P$, $B$, $\varepsilon$
    \item Profiles: $T(z,t)$, $S(z,t)$, $\rho(z,t)$
    \item Scalar timeseries: SST, MLD, surface fluxes
\end{itemize}

Saved to NetCDF files for post-processing and visualization.

\section{Testing and Validation}
\label{sec:testing}

\subsection{Unit Test Suite}

Regression tests in \verb|test_staggering.py| validate individual components:

\textbf{Test coverage}:
\begin{enumerate}
    \item \textbf{Diffusion solver}: Verify implicit solver against analytical solution
    \item \textbf{N² computation}: Validate sign convention (Fix 1) with known profile
    \item \textbf{Shear production}: Check production term $P = \kappa_m S^2$ magnitude
    \item \textbf{TKE dissipation}: Verify dissipation $\varepsilon = c_\varepsilon k^{3/2}/\ell$
    \item \textbf{Prandtl number}: Test $\kappa_h / \kappa_m$ ratio
    \item \textbf{Stratification response}: Verify enhanced mixing in unstable regions
    \item \textbf{Combined forcing}: Test interaction of wind, buoyancy forcing
    \item \textbf{Parameter consistency}: Validate YAML loading, default values
\end{enumerate}

\textbf{Status}: 8/8 tests passing as of 2026-07-16.

\subsection{Scenario-Based Validation}

Six realistic scenarios validate full system behavior:

\textbf{1. arctic\_convection}
\begin{itemize}
    \item Strong cooling, weak wind
    \item Tests convective mixing, EOS clamp (Fix 4)
    \item Final SST: GGL90 $-8.5°$C, KPP $-6.0°$C
\end{itemize}

\textbf{2. calm\_baseline}
\begin{itemize}
    \item Weak forcing, tests background mixing
    \item Final SST: GGL90 21.9°C, KPP 21.9°C (nearly identical)
\end{itemize}

\textbf{3. combined\_storm}
\begin{itemize}
    \item Strong wind + cooling, tests entrainment
    \item Realistic North Atlantic profile (Fix 8)
    \item Final SST: GGL90 8.3°C, KPP 7.2°C
\end{itemize}

\textbf{4. heavy\_rain\_freshening}
\begin{itemize}
    \item Surface freshwater flux, tests salinity mixing
    \item Final SST: GGL90 22.0°C, KPP 22.0°C (nearly identical)
\end{itemize}

\textbf{5. hurricane\_wind}
\begin{itemize}
    \item Extreme wind stress (30 m/s), tests mixing length limiter (Fix 3)
    \item Realistic tropical profile (Fix 8)
    \item Final SST: GGL90 13.5°C, KPP 9.2°C (scheme differences expected)
\end{itemize}

\textbf{6. tropical\_heating\_diurnal}
\begin{itemize}
    \item Diurnal cycle, weak wind, tests surface layer response
    \item Final SST: GGL90 26.4°C, KPP 26.4°C (nearly identical)
\end{itemize}

\textbf{Validation criteria}:
\begin{itemize}
    \item Exit code 0 (successful completion)
    \item Physically reasonable SST evolution
    \item No numerical instabilities (NaN, Inf)
    \item Diagnostics within expected ranges
\end{itemize}

\subsection{Comparison to MITgcm Output}

\textbf{Current status}: Qualitative comparison completed; quantitative comparison in progress.

\textbf{Challenges}:
\begin{itemize}
    \item 1D column vs.\ 3D model (boundary conditions differ)
    \item Simplified EOS in Python port
    \item Different numerical precision (Python float64 vs.\ Fortran real*8 with compiler variations)
\end{itemize}

\textbf{Approach}: Extract single-column output from MITgcm run, compare timestep-by-timestep diagnostics (TKE, $\kappa_h$, SST).

\subsection{Known Limitations}

\begin{enumerate}
    \item \textbf{1D only}: No horizontal advection, mixing, or lateral boundary effects
    \item \textbf{No IDEMIX}: Internal wave energy model not implemented
    \item \textbf{No Langmuir}: Wave-driven mixing absent
    \item \textbf{Simplified EOS}: JMD95 without full compressibility terms
    \item \textbf{No grid-scale smoothing}: MITgcm applies horizontal smoothing to prevent noise; not needed in 1D
\end{enumerate}

\section{Usage Examples and Workflows}
\label{sec:usage}

\subsection{Running Single Experiments}

\textbf{Basic usage} (default parameters):
\begin{verbatim}
cd main
python run_experiment_example.py --scheme ggl90
\end{verbatim}

\textbf{With ECCOv4 R4 parameters}:
\begin{verbatim}
python run_experiment_example.py \
    --scheme ggl90 \
    --ggl90-yaml ../configuration_yamls/ggl90_eccov4r4.yaml
\end{verbatim}

\textbf{Compare both schemes}:
\begin{verbatim}
python run_experiment_example.py --scheme both
\end{verbatim}

\textbf{Enable convective adjustment}:
\begin{verbatim}
python run_experiment_example.py \
    --scheme ggl90 \
    --ivdc-kappa 10.0
\end{verbatim}

\subsection{Batch Scenario Execution}

\textbf{Run all scenarios with defaults}:
\begin{verbatim}
cd main
python run_scenarios.py
\end{verbatim}

\textbf{Run all scenarios with ECCOv4 R4}:
\begin{verbatim}
python run_scenarios.py \
    --ggl90-yaml ../configuration_yamls/ggl90_eccov4r4.yaml \
    --ivdc-kappa 10.0
\end{verbatim}

\textbf{Output}: Summary table with final SST for each scenario/scheme combination.

\subsection{Custom Configuration}

\textbf{Create custom parameter file}:
\begin{verbatim}
# custom_ggl90.yaml
GGL90alpha: 15.0       # Intermediate mixing length
GGL90TKEmin: 1.0e-9    # Higher background TKE
GGL90viscMax: 50.0     # Lower viscosity cap
\end{verbatim}

\textbf{Run with custom parameters}:
\begin{verbatim}
python run_experiment_example.py \
    --scheme ggl90 \
    --ggl90-yaml custom_ggl90.yaml
\end{verbatim}

\subsection{Sensitivity Studies}

\textbf{Parameter sweep example} (vary $\alpha$):
\begin{verbatim}
for alpha in 1.0 5.0 10.0 20.0 30.0; do
    cat > alpha_${alpha}.yaml <<EOF
GGL90alpha: $alpha
EOF
    python run_scenarios.py \
        --ggl90-yaml alpha_${alpha}.yaml \
        --output-dir results_alpha_${alpha}
done
\end{verbatim}

Produces separate result directories for each $\alpha$ value, enabling systematic parameter sensitivity analysis.

\subsection{Common Workflows}

\textbf{1. Quick validation after code changes}:
\begin{verbatim}
# Run regression tests
pytest test_staggering.py

# Run single scenario
python run_experiment_example.py --scheme ggl90 --no-plots
\end{verbatim}

\textbf{2. Full validation suite}:
\begin{verbatim}
# All tests + all scenarios
pytest test_staggering.py && python run_scenarios.py
\end{verbatim}

\textbf{3. Scheme comparison}:
\begin{verbatim}
# Run both schemes, generate comparison plots
python run_experiment_example.py --scheme both
# Output: side-by-side T(z,t), SST(t), mixing coefficient profiles
\end{verbatim}

\textbf{4. Parameter tuning}:
\begin{verbatim}
# Create parameter sweep
python scripts/parameter_sweep.py \
    --param GGL90alpha \
    --values 1.0 10.0 30.0 \
    --scenarios hurricane_wind combined_storm

# Analyze results
python scripts/analyze_sweep.py --sweep-dir results/sweep_GGL90alpha
\end{verbatim}

\subsection{Output Files}

Standard output structure:
\begin{verbatim}
simulations/output/
  <scenario_name>_<scheme>/
    diagnostics.nc       # Full 3D (z,t,variable) diagnostics
    timeseries.nc        # Scalar timeseries (SST, MLD, fluxes)
    profiles.png         # T(z), S(z) initial/final comparison
    timeseries.png       # SST(t), MLD(t) evolution
    mixing_coeff.png     # kappa_h(z,t) heatmap
\end{verbatim}

\subsection{Command-Line Help}

Both scripts provide comprehensive help:
\begin{verbatim}
python run_experiment_example.py --help
python run_scenarios.py --help
\end{verbatim}

\section*{References}

\begin{itemize}
    \item Gaspar, P., Y. Gregoris, and J.-M. Lefevre (1990). A simple eddy kinetic energy model for simulations of the oceanic vertical mixing: Tests at Station Papa and Long-Term Upper Ocean Study site. \textit{Journal of Geophysical Research}, 95(C9), 16,179-16,193.
    \item MITgcm Group (2024). MITgcm User Manual. \url{https://mitgcm.readthedocs.io/}
    \item Jackett, D. R., and T. J. McDougall (1995). Minimal adjustment of hydrographic profiles to achieve static stability. \textit{Journal of Atmospheric and Oceanic Technology}, 12(4), 381-389.
\end{itemize}

\end{document}
